\documentclass[12pt]{amsart}
\usepackage[all]{xy} % for commutative diagrams
\usepackage{amsmath} % various goodies
\usepackage{amsthm} % some stuff for theorem numbering
\usepackage{latexsym}
\usepackage{amsfonts}
\usepackage{amssymb}

%% Theorem-type environments.
% The first parameter to \newtheorem is what the environment is called in the 
% LaTeX source. (What goes into the braces after \begin and \end.)
% The second is the name that shows up in the text.

\theoremstyle{plain} 
\newtheorem{theorem}{Theorem}
\newtheorem{lemma}{Lemma}
\newtheorem{corollary}{Corollary}
\newtheorem{proposition}{Proposition}
\newtheorem{conjecture}{Conjecture}
% (Criterion, Algorithm)

\theoremstyle{definition} 
\newtheorem{definition}{Definition}
\newtheorem{example}{Example}
\newtheorem{todo}{To Do}
% (Condition, Problem)

\theoremstyle{remark} 
\newtheorem{remark}{Remark}
\newtheorem{note}{Note}
% (Notation, Claim, Summary, Acknowledgement, Case, Conclusion)

%% Math operators (easy thanks to amsmath)
\DeclareMathOperator{\Aut}{Aut}
\DeclareMathOperator{\core}{Core}
\DeclareMathOperator{\orb}{Orb}
\DeclareMathOperator{\cay}{Cay}
\DeclareMathOperator{\sab}{Sab}
\DeclareMathOperator{\diag}{diag}
\DeclareMathOperator{\arc}{Arc}
\DeclareMathOperator{\grp}{Grp}
\DeclareMathOperator{\ext}{Ext}

% Commands:
\newcommand{\cB}{\mathcal{B}}
\newcommand{\cS}{\mathcal{S}}
\newcommand{\cP}{\mathcal{P}}
\newcommand{\cZ}{\mathcal{Z}}
\newcommand{\cD}{\mathcal{D}}

% lazy
\newcommand{\arbitraryset}{\mathcal{X}}
\newcommand{\open}{\mbox{``}}
\newcommand{\close}{\mbox{''}}
\newcommand{\wt}[1]{\widetilde{#1}}

\begin{document}

\title{Quotients and Extensions}
\maketitle


%Butterflies in bloom. See how the grow from a small seed into such majestic beings?

\section{Quotient Graphs}


In the first section of this Chapter,
we determined, for any two transitive permutation groups
$(G, \Omega_1)$ and $(G,\Omega_2)$,
the conditions under which there exists a
$G$-homomorphism from $(G,\Omega_1)$ to $(G,\Omega_2)$.

We saw that if $(G,\Omega_1)$
is permutation equivalent to $(G,\cos_G(H))$
for some $H < G$,
and $(G,\Omega_2)$ is permutation equivalent
to $(G,\cos_G(K))$ for some $K < G$,
then there is a $G$-homomorphism from $(G,\Omega_1)$
to $(G,\Omega_2)$ if and only if $H < K < G$.

In this section we give analogous conditions
under which there exists a $G$-homomorphism
from a $G$-symmetric graph $\Gamma$
to a $G$-symmetric graph $\Sigma$.

\begin{theorem} [Lorimer]
If $\Gamma$ is a $G$-symmetric graph isomorphic to $\sab(G,H,HaH)$
and $\Sigma$ is a quotient of $\Gamma$,
then $\Sigma$ is isomorphic to $\sab(G,K,KaK)$ for some $H < K < G$
with $a \not\in K$.
\begin{proof}
Suppose that $\pi : \Gamma \to \Sigma$ is a $G$-graph homomorphism.
Since $\Gamma$ is $G$-isomorphic to $\sab(G,H,HaH)$,
the action of $G$ on the vertices of $\Gamma$ is permutation
equivalent to $(G,\cos_G(H))$.
Since $\pi$ induces a homomorphism from $V \Gamma$ to $V \Sigma$,
by [where?] $(G,V \Sigma)$
must be permutation equivalent to $(G,\cos_G(K))$
for some $K$ with $H < K < G$.
What about the $a$?
\end{proof}
\end{theorem}

\section{Extension Problem}

For a given symmetric graph $\Gamma$, let $\grp(\Gamma)$ denote
the set of subgroups of $\Aut(\Gamma)$ 
which act symmetrically on $\Gamma$.
For a given group $G$, let $\ext(G)$ denote the set of groups $\widetilde{G}$
such that $\widetilde{G}/N \cong G$ for some $N$ normal in $\widetilde{G}$.
Clearly $\Gamma$ is $\widetilde{G}$-symmetric if and only if
$\widetilde{G} \in \ext(G)$ for some $G \in \grp(\Gamma)$.

\begin{definition}
The symmetric graph $(\widetilde{G},\widetilde{\Gamma})$
is said to be an {\it extension} of the symmetric graph $(G,\Gamma)$
if $\widetilde{G} \in \ext(G)$
and there exists a $\widetilde{G}$-homomorphism
from $\widetilde{\Gamma}$ to $\Gamma$.
\end{definition}

The ``extension problem'' for symmetric graph
is the problem of finding all the extensions
of a given symmetric graph $(G,\Gamma)$.

On a group theoretic level, if we knew the subgroup structure of $G$
and we knew all the groups $\widetilde{G}$ which admit $G$ as a 
quotient then in a sense we would ``know'' all the possible 
extensions of $(G,\Gamma)$.

\begin{proposition}
If $\Gamma$ is $G$-isomorphic to $\sab(G,H,HaH)$
then the faithful extensions of $\Gamma$
are the graphs of the form $\sab(G,K,KaK)$
where $K < H$ and $KaK \subset HaH$.
\end{proposition}

\begin{proposition}
For any $\widetilde{G} \in \ext(G)$, let $N \triangleleft G$
be such that $\widetilde{G}/N \cong G$
and let $\pi : \widetilde{G} \to G$ be the natural projection.
The $\widetilde{G}$-symmetric extensions of $\Gamma$
are the graphs of the form 
$\sab(\widetilde{G},R,R\widetilde{a}R))$
where $\pi(R) < H$ and $\pi(a) \in HaH$.
\end{proposition}

From a combinatorial perspective, 
if $(\widetilde{G},\widetilde{\Gamma})$ 
is an extension of $(G,\Gamma)$
then we would like to understand
how the structure of $\widetilde{\Gamma}$
is related to the structure of $\Gamma$.
Ideally, we'd like to be able to describe the adjacency
structure of $\widetilde{\Gamma}$ in terms of the
adjacency structure of $\Gamma$.

\section{Unfaithful Extensions}

Recall from Chapter 1, that in our definition of a symmetric
graph, we do not require the action of the group
to be faithful on the vertices of the graph.
Suppose that $\Gamma$ is a $G$-symmetric graph
isomorphic to $\sab(G,H,HaH)$.
The vertices of $\Gamma$ are the cosets of $H$ in $G$
and the action of $G$ on the vertices of $\Gamma$
is the permutation group $(G,\cos_G(H))$.
By proposition [?] the {\it kernel} of this action is:
$$\core_G(H) = \bigcap_{g \in G} g^{-1}Hg.$$
If this kernel is non-trivial then the action of $G$
on the vertices of $\Gamma$ is unfaithful.

\begin{proposition}
% careful with the a
For any pair of groups $\wt{G}$ and $\wt{H}$ with $\wt{H} < \wt{G}$,
and any $a \in \wt{G}$ with $\wt{a}^2 = 1$ and $\wt{a} \not\in \wt{H}$,
the graphs
$\sab(\wt{G},\wt{H},\wt{H}a\wt{H})$ and $\sab(G, H, HaH)$
are isomorphic where: 
\begin{eqnarray*}
N  & = &\core_{\wt{G}}(\wt{H}) \\ 
G & = & \wt{G}/N \\
H & = & \wt{H}/N
\end{eqnarray*}
and $a = \pi(\wt{a})$ where $\pi : \wt{G} \to G$ is the natural projection.
\begin{proof}
Let $\wt{\Gamma} = \sab(\wt{G},\wt{H},\wt{H}\wt{a}\wt{H})$ 
and let $\Gamma = \sab(G,H,HaH)$.
The vertices of $\wt{\Gamma}$ are the right cosets of $H$ in $G$.
The vertices of $\Gamma$ are the right cosets of $H/N$ in $G/N$.
Now,
\begin{eqnarray*}
[\wt{G}/N : \wt{H}/N] & = & (|\wt{G}|/|N|)/(|\wt{H}|/|N|) \\
            & = & |\wt{G}|/|\wt{H}| \\
	    & = & [\wt{G}:\wt{H}]
\end{eqnarray*}
so both $\wt{\Gamma}$ and $\Gamma$ have the same number of vertices.
Let $\pi : \wt{G} \to G$ be the natural projection.
Since $N$ is normal in $\wt{H}$, we have $\pi(\wt{H}) = H$.
Now,
\begin{eqnarray*}
H\pi(x) = H\pi(y) & \Leftrightarrow & \pi(xy^{-1}) \in H \\
                    & \Leftrightarrow & xy^{-1} \in \wt{H} \\
		    & \Leftrightarrow & \wt{H}x = \wt{H}y \\
\end{eqnarray*}
so $\pi$ actually induces a bijection between the cosets of $H$
in $G$ and the cosets of $H'$ in $G'$.
To see that it preserves adjacency,
suppose that $(\wt{H}x,\wt{H}y)$ is any arc of $\wt{\Gamma}$,
and so $xy^{-1} \in \wt{H}\wt{a}\wt{H}$.
It follows that $\pi(xy^{-1}) \in HaH$
and so $(Hx,Hy)$ is an arc of $\Gamma$.
The result follows.
\end{proof}
\end{proposition}

Suppose that $\wt{\Gamma}$ is a symmetric graph 
isomorphic to $\sab(\wt{G},\wt{H},\wt{H}\wt{a}\wt{H})$.
Suppose further that $\wt{G}$ acts faithfully on the vertices of $\wt{\Gamma}$.
Let $\Gamma$ be any quotient of $\wt{\Gamma}$,
so that, by proposition [?], $\Gamma$ is isomorphic to 
$\sab(\wt{G},\wt{K},\wt{K}\wt{a}\wt{K})$ for some $\wt{K}$
such that $\wt{H} < \wt{K} < \wt{G}$ and $\wt{a} \not\in \wt{K}$.
Although $\wt{G}$ acts faithfully on the vertices of $\wt{\Gamma}$
it does not follow that $\wt{G}$ acts faithfully 
on the vertices of the quotient $\Gamma$.

\begin{definition}
If $(G,\Gamma)$ is a symmetric graph with $G$ acting faithfully on the vertices
of the $\Gamma$, then an {\it unfaithful} extension of $(G,\Gamma)$
is an extension of the form $(\widetilde{G},\widetilde{\Gamma})$
where $\widetilde{G} \neq G$. 
\end{definition}


\begin{proposition}
Suppose that $(G,\Gamma)$ is a symmetric graph with
$\Gamma \cong \sab(G,H,HaH)$.
and that $(\wt{G},\wt{\Gamma})$ is an unfaithful extension
of $(G,\Gamma)$ with
$\wt{\Gamma} \cong 
\sab(\wt{G},\wt{H},\wt{H}\wt{a}\wt{H}))$
If $\pi : \wt{G} \to G$ is the natural projection, then
$\pi(\wt{H}) < H$ and $\pi(\wt{a}) \in HaH$.
\begin{proof}
Let $N = \ker(\pi))$
Since $\Gamma$ is a quotient of $\wt{Gamma}$
we must have 
$\Gamma \cong \sab(\wt{G},\wt{K},\wt{K}\wt{a}\wt{K}))$
for some $\wt{H} < K < \wt{G}$.
Since $\Gamma \cong \sab(G,H,HaH)$
we must have $\core_{\wt{G}}(K) = N$
and $\pi(K) = H$.
Since $\wt{H} < K$ it follows that $\pi(\wt{H}) < \pi(K) < H$.
What about the $a$?
\end{proof}

\end{proposition}

\section{Covering Graphs}

\begin{definition}
The graph $\widetilde{\Gamma}$ is said to be a {\it cover}
of the graph $\Gamma$ if there exists a graph homomorphism
$\pi : \widetilde{\Gamma} \to \Gamma$ with the property that.
For any arc $(B,C)$ of $\Gamma$, and any $v \in \pi^{-1}(B)$
there is a unique $w \in \pi^{-1}(C)$ such that $(v,w)$
is an arc of $\widetilde{\Gamma}$.
\end{definition}

\begin{definition}
The graph $\widetilde{\Gamma}$ is said to be a {\it cover}
of the graph $\Gamma$ if there exists a graph homomorphism
$\pi : \widetilde{\Gamma} \to \Gamma$ with the property that.
For any arc $(B,C)$ of $\Gamma$, and any $v \in \pi^{-1}(B)$
there exists some $w \in \pi^{-1}(C)$ 
(not necessarily unique) such that $(v,w)$
\end{definition}

Clearly a multicover is a special case of a cover.
We shall see that in this section
that if $(\widetilde{G},\widetilde{\Gamma})$
is an unfaithful extension of $(G,\Gamma)$ then
$\widetilde{\Gamma}$ is a multicover of of $\Gamma$.

% Nice geometric description + picture for Cale.

\begin{proposition}
Let $(\widetilde{G},\widetilde{\Gamma})$ be an unfaithful
extension of $(G, \Gamma)$ and let $N$ be a normal subgroup of 
$\widetilde{G}$
such that $\widetilde{G}/N \cong G$.
If $B$ is any vertex of $\Gamma$ then $N$ acts transitively on
the fiber $\pi^{-1}(B)$.
\begin{proof}
Let $\Delta = \pi^{-1}(B)$.
By proposition [?], $\Delta$ is a block of imprimitivity
of $V \widetilde{\Gamma}$.
Suppose that $\Gamma \cong \sab(G,H,HaH)$ and
$\widetilde{\Gamma} \cong
\sab(\widetilde{G},\widetilde{H},\widetilde{H}\widetilde{a}\widetilde{H}))$.
Then the stabilizer of $\Delta$ is $HN$ and the action
of $HN$ on the fiber is permutation equivalent to $(HN,\cos(H))$.

\end{proof}
\end{proposition}


\newpage

\section{Combinatorial Perspective}



In this chapter we look at a number of well-known
``combinatorial'' methods for constructing
extensions of symmetric graphs. We first describe the constructions,
and then analyze them group theoretically.
In the next chapter we describe a ``geometrical'' approach
to the extension problem for symmetric graphs which was
proposed by Gardiner and Praeger.

\section{Direct Sum and Lexicographic Product}

There are a number of ways of forming ``products''
of graphs. See for example [?].

\begin{definition}
For any two graphs $\Gamma$ and $\Sigma$,
then the direct sum of $\Gamma$ and $\Sigma$
is the graph with vertex set $V \Gamma \times V \Sigma$
where $(v,x)$ is adjacent to $(w,y)$ if and only
if $v$ is adjacent to $w$ in $\Gamma$
and $x$ is adjacent to $y$ in $\Sigma$.
\end{definition}

\begin{proposition}
If $\Gamma$ and $\Sigma$ are both symmetric graphs
then the direct product of $\Gamma$ and $\Sigma$
is symmetric.
\end{proposition}

\begin{definition}
For any two graphs $\Gamma$ and $\Sigma$,
the lexicographic product of $\Gamma$ by $\Sigma$
is the graph with vertex set $V \Gamma \times V \Sigma$
where $(v,x)$ is adjacent to $(w,y)$
if and only if $v$ is adjacent to $w$ in $\Gamma$
or $v = w$ and $x$ is adjacent to $y$ in $\Sigma$
\end{definition}

\begin{proposition}
If $\Gamma$ is any symmetric graph and $\Sigma$ is the empty
graph on $n$ vertices, then the lexicographic product
of $\Gamma$ by $\Sigma$ is symmetric.
\end{proposition}

%What would still be lacking is a {\it combinatorial} understanding
%of how the structure of the new graph is related to that of the original.
%On a practical level we would like to have some method of constructiong 
%$\widetilde{\Gamma}$ from $\Gamma$. That is, we'd like to be able to describe
%the adjacency structure of $\widetilde{\Gamma}$ in terms of the adjacency
%structure of $\Gamma$.

%It may be the case that no uniform description exists for all pairs
%$(\Gamma, \widetilde{\Gamma})$ where $\Gamma$ is a quotient of $\widetilde{\Gamma}$.
%Nevertheless it should be possible to identify families of pairs which admit
%some nice description. These families will be characterized as satisfying
%some additional combinatorial or group theoretic property. The task is to find
%which additional constraints will lead to a tight combinatorial coupling between 
%the symmetric graph $\Gamma$ and its ``extension'' $\widetilde{\Gamma}$.
%
%Again there is no reason to believe that any one construction should be suitable
%for all possible pairs $(\Gamma, \widetilde{\Gamma})$.
%It would be interesting however to gather together those constructions which are known
%and analyze the {\it group theoretic} properties which the pairs 
%$(\Gamma, \widetilde{\Gamma})$ related by this {\it combinatorial} construction share.

%\begin{todo}
%Talk about lexicographic products and direct sums of graphs.
%Talk about Bigg's covering graph construction.
%\end{todo}


%\section{This goes elsewhere}
%We must pause for a moment to highlight an important distinction
%between a symmetric graph and a $G$-symmetric graph.
%We say that a graph $\Gamma$ is symmetric if it is $\Aut(\Gamma)$-symmetric.
%Simmilarly we say that a design $\cD$ is flag-transitive if it is
%$\Aut(\cD)$-flag transitive.

%It is often the case however that for a given symmetric graph $\Gamma$
%there is a subgroup $G$ of the full automorphism group
%which also acts symmetrically on the graph.
%In this case the graph will be $G$-symmetric as well as $\Aut(\Gamma)$-symmetric.
%We saw an example of this in Chapter 1 when we constructed the cube
%as an orbital graph of $(S_4, \cos(C_3))$ but found its full automorphism group
%is in fact $S_4 \times Z_2$.
%A simmilar ``problem'', also arises with flag-transitive designs.
%
%
%on $\Gamma$. Simmilarly, for a given flag-transitive design, $\cD$
%let $\grp(\cD)$ denote the set of subgroups (upto isomorphism)
%of $\Aut(\cD)$ which act symmetrically on $\cD$.
%Then, for a given group $G$, $\Gamma$ is $G$-symmetric
%if and only if $G$ admits some element of $\grp(\Gamma)$ as a composition factor.
%Simmilarly $\cD$ is $G$-flag transitive if and only if
%$G$ admits some element of $\grp(\cD)$ as a composition factor.
%
%The study of $G$-symmetric graphs can thus be approached in two different ways.
%The first is to fix the group $G$ and examine the family of graphs $\Gamma$
%which admit a symmetric action of $G$, in the hope that this will reveal
%some information about the group $G$.
%The second approach is to fix a graph $\Gamma$ and study the groups
%which act symmetrically on $\Gamma$ in the hope that this will
%reveals some information about the graph $\Gamma$.


\section{Bigg's Covering Graph Construction}

% Introductory paragraph:
The three-arc graph construction allows us, under certain conditions,
to ``unfold'' a $G$-symmetric graph into a larger, {\it imprimitive}
$G$-symmetric graph, admitting the original graph as a quotient.
The {\it covering graph} construction \cite{B1974} is simmilar
in this respect. It also, under certain conditions,
allows us to ``unfold'' a $G$-symmetric graph into a larger one.
It differs from the three-arc graph construction, however,
in that it requires a simultaneous ``unfolding''
of the group $G$.

% Semidirect Product:
Recall that given two groups $N$ and $G$ and a homomorphism
$$\rho : G \to \Aut(N)$$
we may form the {\it semidirect product} of $N$ by $G$.
This is the group $\widetilde{G} = N \rtimes_\rho G$, 
whose elements are the ordered pairs: 
$$\{ (n,g) : n \in N, g \in G \}$$
with multiplication given by:
$$(n_1, g_1)(n_2,g_2) = (n_1^{\rho(g_1)}n_2, g_1g_2).$$
The functions $i_1 : N \to \widetilde{G}$ and $i_2 : G \to \widetilde{G}$
given by $ n \mapsto (n,1) $ and $ g \mapsto (1,g) $ respectively give
natural embeddings of $N$ and $G$ into $\widetilde{G}$.
Identifying $N$ with its image under $i_1$,
the semidirect product has the property that $N$ is normal in $\widetilde{G}$
and $\widetilde{G}/N \cong G$.

% Definitions:
\begin{definition} [$N$-chain]
Suppose that $\Gamma$ is a $G$-symmetric graph and $N$ is a group.
An $N$-chain is a function $\phi : A \Gamma \to N$ 
satisfying $\phi((v,u)) = \phi((u,v))^{-1}$ for all $(u,v) \in A \Gamma$.
\end{definition}

\begin{definition} [Compatible $N$-chain]
Suppose that $\Gamma$ is a $G$-symmetric graph, $N$ is a group
and $\rho : G \to \Aut(N)$ is a homomorphism.
An $N$-chain $\phi$ is said to be {\it compatible} with $\rho$
if for every $g \in G$ the following diagram commutes:

\centerline{
\xymatrix{
\ar[d]_{g} \ar[r]^\phi A \Sigma & K \ar[d]^{\rho(g)} \\
\ar[r]_{\phi} A \Sigma & K
}}
\end{definition}

\begin{definition}[Biggs Cover]
Suppose that $\Gamma$ is a $G$-symmetric graph, $N$ is a group,
$\rho : G \to \Aut(N)$ is a homomorphism and $\phi$ is a compatible $N$-chain.
The Biggs Cover $\widetilde{\Gamma}(N,\rho,\phi)$
of $\Gamma$ with respect to $\phi$ is the graph 
with vertex set: 
$\{ (g,v) : k \in G, v \in V \Gamma \}$
and arc set: 
$ \{ ((g_1,v_1),(g_2,v_2)) : (v_1,v_2) \in A \Gamma, g_2 = g_1 \phi((v_1,v_2)) \}$
\end{definition}

Note that the condition $\phi((v,u)) = \phi((u,v))^{-1}$ ensures that the resulting 
graph is simple.

% Proposition:
\begin{proposition}
The covering graph $\widetilde{\Gamma}(N,\rho,\phi)$ is $(N \rtimes_\rho G)$-symmetric.
\begin{proof}

Define an action of $N \rtimes_\rho G$ on $\widetilde{\Gamma}$ by:
$$(n,v)^{(\eta, g)} = (n^{\rho(g)} \eta,v^g).$$
This action is well-defined since:
\begin{eqnarray*}
(n,v)^{(\eta_1, g_1)(\eta_2, g_2)} 
& = & (n^{\rho(g_1)} \eta_1 ,v^{g_1})^{(\eta_2, g_2)} \\
& = & (n^{\rho(g_1) \rho(g_2)} \eta_1^{\rho(g_2)} \eta_2 , v^{g_1 g_2}) \\
& = & (n,v)^{(\eta_1^{\rho(g_2)} \eta_2, g_1 g_2)}
\end{eqnarray*}
We must check that it preserves the adjacency structure of $\widetilde{\Gamma}$.
Observe that, by the compatibility of $\rho$ and $\phi$, 
if $n_2 = n_1 \phi(v_1,v_2)$ then:
\begin{eqnarray*}
{n_2}^{\rho(g)} \phi({v_1}^g, {v_2}^g)
& = & {n_2}^{\rho(g)} \phi(v_1, v_2)^{\phi(g)} \\
& = & (n_2 \phi(v_1, v_2))^{\phi(g)} \\
& = & {n_1}^{\phi(g)} 
\end{eqnarray*}
for any $(\eta,g) \in N \rtimes_\rho G$.
Also, since $G$ is a group of automorphisms of $\Gamma$, we know that
$(v_1,v_2) \in A \Gamma$ if and only if $({v_1}^g,{v_2}^g) \in A \Gamma$.
Thus the action defined above does indeed preserves adjacency.

\hfill \break
For any $(n_1,v_1), (n_2,v_2) \in V \widetilde{\Gamma}$, 
by the $G$-symmetry of $\Gamma$ we can find a $g \in G$ 
such that ${v_1}^g = {v_2}$.
Let $\eta = {n_1}^{\rho(g)^{-1}} n_2$.
We have:
\begin{eqnarray*}
(n_1,v_1)^{(\eta,g)} & = & ({n_1}^{\rho(g)} \eta,{v_1}^g) \\
                    & = & ( {n_1}^{\rho(g)}, {n_1}^{\rho(g)^{-1}} n_2,v_2) \\
		    & = & (n_2,v_2).
\end{eqnarray*}
Thus the action is vertex transitive.

\hfill \break
Suppose that $(n_1,v_1)$ and $(n_2,v_2)$ are both adjacent to $(n,g)$,
so $n_1 = n \phi(v,v_1)$ and $n_2 = n \phi(v,v_2)$.
By the $G$-symmetry of $\Gamma$ we can find a $g \in G_v$ such that
$v_1^g = v_2$.
Let $\eta = (n^{-1})^{\rho(g)}n$. We have:
\begin{eqnarray*}
(n,v)^{(\eta,g)} & = & (n^{\rho(g)} \eta, v^g) \\
                 & = & (n^{\rho(g)} (n^{-1})^{\rho(g)}n, v) \\
		 & = & (n,v)
\end{eqnarray*}
and:
\begin{eqnarray*}
(n_1,v_1)^{(\eta,g)} & = & ({n_1}^{\rho(g)} \eta, {v_1}^g) \\
                 & = & ({n_1}^{\rho(g)} (n^{-1})^{\rho(g)}n, v_2) \\
                 & = & ({(n_1 n^{-1})}^{\rho(g)} n, v_2) \\
                 & = & (\phi(v,v_1)^{\rho(g)} n, v_2) \\
                 & = & (\phi(v,v_2) n, v_2) \\
                 & = & (n_2 n^{-1} n, v_2) \\
                 & = & (n_2, v_2) \\
\end{eqnarray*}
Thus the action is locally transitive.
The result follows.

\end{proof}
\end{proposition}

In fact, a stronger result than this holds.
If $\Gamma$ is $(G,s)$-arc transitive
then $\widetilde{\Gamma}$ is $(\widetilde{G},s)$-arc transitive,
where $\widetilde{G} = N \rtimes_\rho G$.
The proof is by induction and may be found in \cite{B1974}.
This construction was originally used by Conway
to produce an infinitely family of $5$-arc transitive graphs.

\hfill \break
For each $v \in V \Gamma$, let $B(v) = \{ (n,v) : n \in N \}$.
For any $(\eta, g) \in \widetilde{G}$ we have:
\begin{eqnarray*}
B(v)^{(\eta,g)} & = & \{ (n,v)^{(\eta,g)} : n \in N \} \\
                & = & \{ (\eta n^{\rho(g)},v^g) : n \in N \} \\
		& = & \{ (n, v^g) : n \in N \} \\
		& = & B(v^g).
\end{eqnarray*}
Thus the partition: 
$\cB = \{ B(v) : v \in V \Gamma \}$
is $\widetilde{G}$-invariant.

\begin{proposition}
The quotient of the Biggs cover $\widetilde{\Gamma}$
with respect to the partition $\cB$
is isomorphic to the original graph $\Gamma$
\begin{proof}
The map $\eta : V \Gamma \to V \widetilde{\Gamma}_\cB$ given by
$v \mapsto B(v)$ establishes a bijection between the vertices of $\Gamma$
and the vertices of $\widetilde{\Gamma}_\cB$.
If $(u,v)$ is an arc of $\Gamma$, then $((1,u),(\phi((u,v)),v))$
is an arc of $\widetilde{\Gamma}$, and so $(B(u), B(v))$ 
is an arc of $\widetilde{\Gamma}_\cB$.
If $(u,v)$ is not an arc of $\Gamma$, then for all $n_1, n_2 \in N$
$((n_1, u),(n_2, v))$ is not an arc of $\widetilde{\Gamma}$
and so $(B(u), B(v))$ is not an arc of $\widetilde{\Gamma}_\cB$.
Thus $\eta$ establishes an isomorphism between
$\Gamma$ and $\widetilde{\Gamma}_\cB$.

\end{proof}
\end{proposition}

\newpage
Let $(B,C)$ be any arc of $\widetilde{\Gamma}$.
By the above $B = B(u)$ and $C = B(v)$ for some $(u,v) \in A \Gamma$.
Its not too hard to see that each $(n,u) \in B$ 
has a unique neighbour in $C$, namely $(n \phi(u,v),v)$.
Thus the induced bipartite graph $\Gamma[B,C]$ is a matching.
It follows immediately that the valency of $\widetilde{\Gamma}$
is the same as the valency of $\Gamma$. 

\section{Group Theoretic Analyais of Bigg's Covering Graph Construction}

The group $\widetilde{G}$ acts unfaithfully on $\cB$ with kernel 
$$\widetilde{N} = \{ (n,1) : n \in N \} \cong N.$$
If $|N| = n$ then $|\widetilde{G}| = n |G|$,
and since $\widetilde{\Gamma}$ has exactly $n$ times as many vertices as $\Gamma$, 
we have $[\widetilde{G} : \widetilde{G}_{B(\alpha)}] = n [G : G_\alpha]$.
It follows that 
$ |\widetilde{G}_{B(\alpha)}| = |G_\alpha|$.


Let $\pi : \widetilde{G} \to G$ be the natural projection.
Clearly $|\pi(\widetilde{G}_{B(\alpha)})| \leq |\widetilde{G}_{B(\alpha)}|$.
Since $G_\alpha \leq \pi(\widetilde{G}_{B(\alpha)})$,
we must have $G_\alpha = \pi(\widetilde{G}_{B(\alpha)})$
and $|\pi(\widetilde{G}_{B(\alpha)})| = |\widetilde{G}_{B(\alpha)}|$, 
so $\widetilde{G}_{B(\alpha)} \cong \pi(\widetilde{G}_{B(\alpha)}) \cong G_\alpha$.
In fact, the permutation groups 
$(G_\alpha, \Gamma(\alpha))$ and 
$(\widetilde{G}_{B(\alpha)}, \widetilde{\Gamma}(B(\alpha)))$
are equivalent via the bijection 
$\eta : \Gamma(\alpha) \to \widetilde{\Gamma}(B(\alpha))$
which sends $\beta$ to $B(\beta)$.
Probably a much easier way to say this.

If $\Gamma$ is isomorphic to the coset graph $\Gamma(G, H, a)$
then since the local actions of $\Gamma$ and $\widetilde{\Gamma}$
are the same, $\widetilde{\Gamma}$ must be isomorphic to the coset graph
$\Gamma(\widetilde{G}, H, \widetilde{a})$ where $\widetilde{a}$
is such that $\pi(\widetilde{a}) = a$. That is 
$\widetilde{a} \in \{ (n,a) : n \in N \}$.
Thus for any $N$ there are exactly $n$ possible Biggs covers of $\Gamma$.

Looking closely at the compatibility condition between
$\rho$ and $\phi$ we see that, since $G$ acts transitively
on the arcs of $\Gamma$, the map $\phi$ is completely
determined by where it sends the arc $(v_{H}, v_{Ha})$.
Choosing where to send this arcs is essentially equivalent to
choosing which preimage of $a$ we want.
Requiring that $\phi((v_{H}, v_{Ha})) = n$ is essentially
the same as determining that $\widetilde{a} = (n,a)$.

I think that Bigg's covering graph construction encompasses
all pairs of graphs $\Gamma(G, H, a)$ and $\Gamma(\widetilde{G}, H, \widetilde{a})$
where $\widetilde{G}$ is a semidirect product of $N$ by $G$ for some $N$.
In the more general case where $\widetilde{G}$ is an extention of $N$ by $G$
but not a split extension $\Gamma(\widetilde{G}, H, \widetilde{a})$
should be a cover of $\Gamma(G, H, a)$ in the sense that the induced
bipartite graph is a matching, but Bigg's covering graph construction
cannot be used to construct  $\Gamma(\widetilde{G}, H, \widetilde{a})$
from $\Gamma(G, H, a)$.

\section{Generalization of Bigg's Covering Graph Construction}

The Bigg's covering graph construction applies to the case
where $\widetilde{G}$ is a semidirect product
of $N$ by $G$ and $H < \widetilde{G}$ is such that $H \cap N = \{1\}$.
The graph $\widetilde{\Gamma}$ is isomorphic to 
$\sab(\widetilde{G},H,HaH)$
and the graph $\Gamma$ is isomorphic
to $\sab(\widetilde{G},\widetilde{H},\widetilde{H}a\widetilde{H})$
where $\widetilde{H} = HN$.

In this section we consider a generalization of the 
Bigg's covering graph construction in which the condition
that $H \cap N = \{1\}$ is relaxed.
We find that in this case $\widetilde{\Gamma}$ is a 
{\it multicover} of $\Gamma$.

\begin{lemma}
If $H$ is a subgroup of $G$ and $N$ is a normal subgroup of $G$
then $H \cap N$ is a normal subgroup of $H$.
\end{lemma}

Let $M = N / H \cap N$ and let $\lambda : N \to M$ be the natural projection.


%This would happen,
%for example, if we could find a group $\hat{H} \leq H$
%such that $[H: a^{-1}Ha \cap H] = [\hat{H}: a^{-1}\hat{H}a \cap \hat{H}]$.
%Then the graph $ \widetilde{\Gamma} = \Gamma(G, \hat{H}, a)$ 
%would be a cover of $\Gamma = \Gamma(G, H, a)$,
%but the local actions are not the same.

%Why is it necessarily a cover? Suppose that $[H, \hat{H}] = k$
%then $\widetilde{\Gamma}$ has $k$ times as many vertices as $\Gamma$.
%If both graphs have valency $r$, then there are $rk$ edges eminating
%from any block. Each block is adjacent to $r$ blocks,
%so there must be $k$ edges between any two blocks.



\bibliography{references}
\bibliographystyle{plain}
\end{document}
